\thispagestyle{plain}

\section*{Kurzfassung}

Es wird ein Nächste-Nachbar Tight-Binding-Modell mit BCS-Wechselwirkung zur Untersuchung der Supraleitung auf dem Bienenwabengitter verwendet.
Nach Herleitung der Selbstkonsistenzgleichung für den Ordnungsparameter $\Delta$ wird analytisch gezeigt, dass bei halber Füllung ein quantenkritischer Punkt auftritt.
Die kritische Wechselwirkung $U_C$ ist dabei invers proportional zur Debye-Energie.
Die Abhängigkeiten des Ordnungsparameters und der kritischen Temperatur von verschiedenen Parametern, insbesondere des chemischen Potentials, werden dargestellt.
Das Modell wird auf Graphen angewandt. 
Eine einfache Abschätzung zeigt, dass konventionelle Supraleitung bei halber Füllung ausgeschlossen werden kann, während Graphen bei Dotierung bis zur Van-Hove-Singularität der Zustandsdichte potentiell supraleitende Eigenschaften aufweisen könnte.

\section*{Abstract}
\begin{foreignlanguage}{english}
A nearest-neighbor tight-binding model with BCS interaction is used to study superconductivity on a honeycomb lattice.
After deriving the self-consistency equation for the order parameter $\Delta$, it is shown analytically that a quantum critical point occurs at half filling.
The critical interaction $U_C$ is inversely proportional to the Debye energy.
The dependencies of the order parameter and the critical temperature on various parameters, in particular the chemical potential, are presented.
The model is applied to graphene. 
A simple estimate shows that conventional superconductivity can be ruled out at half-filling, whereas graphene could potentially exhibit superconducting properties when doped up to the Van Hove singularity of the density of states.
\end{foreignlanguage}